\BourbakiBackMatter{Historical Note}

We have seen (historical notes of Chapters I and II--III) that the first noncommutative algebras appear in 1843--44, in works of Hamilton {[}25{]} and Grassmann {[}23{]} and {[}24{]}. Hamilton, when introducing the quaternions, already has a clear idea of arbitrary algebras of finite rank over the field of real numbers ({[}25{]}, Preface, p.~26--31) \(^{(1)}\) . In developing his theory, he later has the idea of considering what he calls ``biquaternions,'' that is, the algebra over the field of complex numbers with the same multiplication table as the quaternion algebra; on this occasion, he observes that this extension causes the appearance of zero divisors ({[}25{]}, p.~650). Grassmann's point of view is slightly different, and for a long time, his ``exterior algebra'' would remain separate from the general theory of algebras; \(^{(2)}\) but in his language, which still lacks precision, one cannot fail to recognize the first idea of an algebra (of finite dimension or not, over the field of real numbers) defined by a generating set and relations {[}24{]}.

New examples of algebras are introduced, more or less explicitly, between 1850 and 1860. Cayley, when developing the theory of matrices {[}8{]}, does not yet see the square matrices as forming an algebra (a point of view that will first be clearly expressed by the Peirces around 1870 {[}40{]}). However, he does already note, on this occasion, the existence of a system of matrices of order 2 satisfying the multiplication table of the quaternions; this observation can be seen as the first example of a linear representation of an algebra. \(^{(3)}\) Moreover, in the memoir where he defines the abstract notion of finite group, he also gives, in passing, the definition of the algebra of such a group, without further applying the definition ({[}7{]}, p.~129).

No other progress is made before 1870, but at this time, research into the general structure of finite-dimensional algebras (over the field of real or complex numbers) begins. It is B. Peirce who takes the first steps in this direction; he introduces the notions of nilpotent element and of idempotent element, proves that an algebra (with or without unit) with at least one nonnilpotent element has a nonzero idempotent, writes down the famous decomposition

\[
x = e x e + (x e - e x e) + (e x - e x e) + (x - x e - e x + e x e)
\]

(e an idempotent, x an arbitrary element), and has the (still slightly imprecise) idea of a decomposition of an idempotent as a sum of pairwise orthogonal ``primitive'' idempotents {[}40{]}. Moreover, according to Clifford ({[}11{]}, p.~274), \(^{(4)}\) it is to B. Pierce that should be attributed the notion of tensor product of two algebras, which Clifford himself applies implicitly to a generalization of Hamilton's ``biquaternions'' {[}9{]} and explicitly to the study of the algebras named after him, several years later ({[}10{]} and {[}11{]}). These new notions are used by B. Peirce to classify low-dimensional algebras (over the field of complex numbers). This problem was also tackled, around 1880, by other mathematicians of the Anglo-American school, in particular Cayley and Sylvester. The great variety of possible structures quickly becomes apparent, and it is undoubtedly this fact that will, in the following period, direct research toward obtaining classes of algebras with more particular properties.

In continental Europe, where the ideas evolve quite differently, such research already appears before 1880. In 1878, Frobenius proves that the quaternions form the only example of a noncommutative (finite-dimensional) field over the field of real numbers ({[}18{]}, p.~59--63), a result published independently two years later by C. S. Pierce {[}39{]}. As early as 1861, Weierstrass, in clarifying a remark by Gauss, had, in his lectures, characterized commutative algebras over R or C without nilpotent elements \(^{(5)}\) as direct sums of fields isomorphic to R or C. Dedekind, in his research, had reached the same conclusions around 1870, in connection with his ``hypercomplex'' conception of the theory of commutative fields. Their proofs are published in 1884--85 ({[}54{]} and {[}12{]}). It is also in 1884 that H. Poincaré, in a quite cryptic note {[}41{]}, draws attention to the possibility of viewing the equations \(z_{i} = \varphi_{i}(x_{1}, \ldots, x_{n}, y_{1}, \ldots, y_{n})\) that express the multiplicative law \((\sum_{i} x_{i} e_{i})(\sum_{i} y_{i} e_{i}) = \sum_{i} z_{i} e_{i}\) in an algebra as defining (locally, of course) a Lie group. This remark seems to have made a great impression on Lie and his followers (E. Study, G. Scheffers, F. Schur, and a bit later T. Molien and É. Cartan) who, at that time, were developing the theory of ``continuous'' groups and especially classification problems (see, in particular, {[}43{]}, p.~387). During the period 1885--1905, it leads the mathematicians of this school to apply methods similar to those they use when studying Lie groups and Lie algebras to the study of the structure of algebras. These methods are based, above all, on the consideration of the characteristic polynomial of an element of the algebra with respect to its regular representation (a polynomial already encountered in the works of Weierstrass and Dedekind mentioned above) and on the decomposition of this polynomial into irreducible factors. That decomposition, as Frobenius would discover a little later, is reflected in the decomposition of the regular representation into irreducible components.

In the course of the research carried out by Lie's school on algebras, the ``intrinsic'' notions of the theory gradually emerge. The notion of radical appears in a specific case (that when the quotient by the radical is a direct product of fields) in G. Scheffers' work in 1891 {[}43{]}, and more clearly in that of T. Molien {[}31{]} and É. Cartan {[}4{]}, who study the general case (the word radical comes from Frobenius {[}22{]}). Study and Scheffers {[}43{]} highlight the concept of an algebra that is a direct product of several others (already glimpsed by Peirce ({[}38{]}, p.~221). Finally, Molien {[}31{]} introduces quotient algebras of an algebra, a notion essentially equivalent to that of two-sided ideal (first defined by É. Cartan {[}4{]}) or of homomorphism (a name also due to Frobenius). The analogy with groups is obvious, and later, in 1904, Epsteen and Wedderburn will consider composition series of two-sided ideals and extend the Jordan--Hölder theorem to them. The most important results of this period are T. Molien's {[}31{]}: guided by the notion of simple group, he defines simple algebras (over C) and proves that these are matrix algebras. He then proves that the structure of an arbitrary algebra of finite rank over C essentially reduces to the case, already studied by Scheffers, when the quotient by the radical is a direct sum of fields. Shortly after that, these results are rediscovered and established more rigorously and more clearly by É. Cartan {[}4{]}. On this occasion, he introduces the notion of semisimple algebra and brings to light numerical invariants, the ``É. Cartan integers'' (VIII, p.~180, Exercise 8), associated with an arbitrary algebra over the field C. He thus brings the theory of these algebras to a point beyond which little progress has been made since. \(^{(6)}\) Finally, he extends Molien's results and his own on algebras over R.

Around 1900, ideas begin to develop in a direction that leads to abandoning every restriction on the field of scalars in everything related to linear algebra. In particular, we must point out the vigorous impulse given to studying finite fields by the American school surrounding E. H. Moore and L. E. Dickson. The result of this research is Wedderburn's theorem {[}50{]} that proves that every finite field is commutative. In 1907, Wedderburn takes Cartan's results and extends them to an arbitrary base field {[}51{]}. In doing so, he completely abandons his predecessors' methods (that cannot be applied when the field is no longer algebraically closed or maximal ordered). Instead, he returns, while refining it, to B. Peirce's technique of idempotents, which allows him to put in definite form the structure theorem for semisimple algebras, reducing their study to that of noncommutative fields. Moreover, the problems of the extension of scalars arise naturally in his perspective, and Wedderburn proves that every semisimple algebra remains semisimple after a separable extension of the base field \(^{(7)}\) and becomes a direct product of central simple matrix algebras when this extension is sufficiently large ({[}51{]}, p.~102). \(^{(8)}\) A bit later, Dickson, for n = 3 {[}17{]}, and Wedderburn himself, for n arbitrary {[}52{]}, give the first examples of noncommutative fields of rank \(n^{2}\) over their center, \(^{(9)}\) thus introducing in a specific case the theory of ``cross products'' and of ``systems of factors'' that would later be developed by R. Brauer {[}2{]} and E. Noether ({[}33{]}, {[}34{]}, and {[}35{]}). Finally, in 1921, Wedderburn proves a specific case of the commutation theorem {[}53{]}.

In the meantime, from 1896 to 1910, a theory close to that of algebras, the theory of linear representations of groups (limited at first to representations of finite groups), was developed by Frobenius, Burnside, and I. Schur. It finds its origin in remarks made by Dedekind. In the course of his research on normal bases of Galois extensions, Dedekind had (even before the publication of his work on algebras), around 1880, come across the ``Gruppendeterminant'' \(\det(x_{st^{-1}})\) , where \((x_{s})_{s\in G}\) is a sequence of variables whose index set is a finite group G (in other words, the norm of the generic element of the algebra of the group G with respect to its regular representation). He had observed that when G is abelian, this polynomial decomposes into linear factors (which generalizes an identity proved long before for determinants of ``circulant'' matrices, which correspond to cyclic groups G). In the course of his fascinating correspondence with Frobenius {[}13{]}, Dedekind, in 1896, draws his attention to this property, its link to the theory of characters of abelian groups, and analogous results on specific noncommutative groups that he obtained in 1886. A few months later, Frobenius completely solves the problem of the decomposition of the ``Gruppendeterminant'' into irreducible factors {[}19{]}, thanks to his brilliant generalization of the notion of character {[}20{]}, which we will not discuss here. However, it should be noted that in the later development of this theory, \(^{(10)}\) Frobenius always remains aware of its connection to the theory of algebras (on which Dedekind had not ceased to insist in his letters). After having introduced the notions of irreducible representation and of completely reducible representation for groups {[}21{]} and proved that the regular representation contains all irreducible representations, it is through analogous methods that he proposes, in 1903, to take up the theory of Molien--Cartan {[}22{]}. In the work of Burnside {[}3{]} and I. Schur {[}45{]}, the ``hypercomplex'' aspect of the theory does not explicitly intervene; but it is in their work that the fundamental properties of irreducible representations, Schur's lemma and Burnside's theorem, appear. Finally, it should be noted that it is in this theory that two particular cases of the commutation theorem appear for the first time. The first is in I. Schur's thesis {[}44{]}, which links (precisely by commutation in the endomorphism ring of a tensor space) the representations of the linear group and those of the symmetric group. The second is in Schur's 1905 work {[}46{]}, where he shows that matrices that commute with all matrices of an irreducible representation over the field C are scalar multiples of \(I\) (a result that also follows from Burnside's theorem).

It remained to clearly identify the common base of these theories: this was the work of the German school surrounding E. Noether and E. Artin, in the period 1925--1933, which sees the creation of modern algebra. Already in 1903, in a memoir on the algebraic integration of linear differential equations {[}42{]}, H. Poincaré had defined, in an algebra, left and right ideals and the notion of minimal ideal. He had also observed that in a semisimple algebra, every left ideal is the direct sum of its intersections with the simple components and that in the algebra of matrices of order n, the minimal ideals have dimension n; but his work went unnoticed by algebraists. \(^{(11)}\) In 1907, Wedderburn defines left and right ideals of an algebra again and proves some of their properties (in particular, that the radical is the largest nilpotent left ideal ({[}51{]}, p.~113--114)).

But it is not until 1927 that these notions are used in an essential way in the theory of algebras. \(^{(12)}\) Generalizing proof methods that had appeared here and there, \(^{(13)}\) W. Krull in 1925 {[}28{]} and E. Noether in 1926 {[}33{]} introduce and systematically use the maximal and minimal conditions. The first uses them to extend Remak's theorem on the decomposition of a finite group as a direct product of indecomposable groups to abelian groups with operators (that he defines on this occasion) (cf.~VIII, p.~37, Theorem 2), while the second uses these conditions in the characterization of Dedekind rings. In 1927, E. Artin {[}1{]}, applying the same idea to noncommutative rings, shows how, through a systematic study of minimal ideals, one can extend Wedderburn's theorem to all rings whose left ideals satisfy both the maximal and the minimal conditions. \(^{(14)}\)

Elsewhere, in 1926, Krull {[}29{]} makes the link between the notion of abelian group with operators and that of linear representation of a group. This point of view is generalized to algebras and developed in detail by E. Noether in a fundamental 1929 work {[}34{]}. Because of the importance of the introduced ideas and the clarity of the exposition, that work deserves to be listed next to Steinitz' memoir on commutative fields as one of the pillars of modern algebra. \(^{(15)}\)

Finally, in a series of works beginning in 1927 ({[}36{]}, {[}2{]}, and {[}35{]}), E. Noether and R. Brauer, joined in 1929 by A. Albert and H. Hasse, take up the study of skew fields where Wedderburn and Dickson left off. Although the most essential part of their results consists in a thorough study of the Brauer group, in particular over algebraic number fields, it is in these works that are clarified the commutation theorems as well as the notion of neutralizing field of a simple algebra and its relation to maximal commutative subfields. Finally, in 1927, Skolem characterizes the automorphisms of simple rings {[}49{]} in a theorem recovered several years later by E. Noether {[}35{]} and R. Brauer {[}2{]}.

Thus, in 1934, the elementary theory of simple and semisimple rings has more or less reached its final form (for an overview of the current state of the theory, see {[}16{]}).
% semantic-fragments: legacy-input-seam
\relax{}
